% MATH companion editorial bridge VI/28; explanatory prose only.
\paragraph{What normalization repairs.}
Let \(A\) be a domain with fraction field \(K\).  Its integral closure
\(\overline A\subseteq K\) consists of all elements of \(K\) satisfying
some monic polynomial with coefficients in \(A\).  The normalization
of \(\operatorname{Spec}A\) is
\[
 \operatorname{Spec}\overline A\longrightarrow\operatorname{Spec}A.
\]
It enlarges the coordinate ring without changing the function field.
Thus normalization is \emph{birational}; in favorable geometric
situations it separates certain branches or repairs singular behavior.
It is not, however, a general procedure for making a scheme smooth.
A normal variety can still have singular points.

\paragraph{A complete example: the cuspidal cubic.}
Assume \(k\) is a field and consider
\[
 A=k[x,y]/(y^2-x^3)\cong k[t^2,t^3]\subseteq k[t].
\]
The parametrization \(x=t^2\), \(y=t^3\) makes
\(t=y/x\) in the common fraction field (where the quotient is
understood as a rational function).
The element \(t\) is integral over \(A\), because it satisfies the
monic equation \(T^2-x=0\).  Consequently
\(A\subseteq k[t]\) is an integral extension.
Conversely \(k[t]\) is integrally closed, since it is a polynomial
ring in one variable over a field.  Thus the integral closure of
\(A\) in \(k(t)\) is exactly \(k[t]\):
\[
 \widetilde X=\operatorname{Spec}k[t]
 \longrightarrow
 X=\operatorname{Spec}k[t^2,t^3],
 \qquad t\longmapsto (t^2,t^3).
\]
The origin \((0,0)\) is the cusp, and the normalization map replaces its
local nonnormal algebra by the regular parameter \(t\).
This example is \emph{finite} because \(k[t]\) is generated by
\(1,t\) as an \(A\)-module: higher powers reduce using \(t^2=x\).

\paragraph{Why localization is the gluing mechanism.}
Normalization must be defined consistently on overlapping affine
charts.  For a multiplicative set \(S\subset A\), integral closure
commutes with localization:
\[
 \overline{S^{-1}A}=S^{-1}\overline A
 \quad\text{inside }\operatorname{Frac}(A).
\]
The forward inclusion follows because a localized integral element
still satisfies a monic polynomial.  For the converse, clear the
finitely many denominators in the monic relation for an element
integral over \(S^{-1}A\).  After multiplying the element by a
suitable denominator \(s\in S\), that element is integral over
\(A\), hence belongs to \(\overline A\).
This fact explains why the affine normalizations agree on every
principal-open overlap, and hence why they glue to a scheme.

\paragraph{A hypothesis that must not be suppressed.}
The normalization of an integral scheme exists without assuming that
the normalization morphism is finite.  Finiteness is an additional
property and does not hold for all integral schemes or all domains.
It does hold in many standard textbook settings: for example, an
integral scheme of finite type over a field has a finite normalization
morphism.  Any proof that uses finiteness of the integral closure as a
module must therefore state the appropriate finiteness hypothesis.

\paragraph{How the universal property is used.}
Suppose \(Z\) is normal and integral and \(f:Z\to X\) is dominant,
where \(X\) is integral.  On a compatible affine chart, elements
integral over \(\mathcal O_X\) pull back to elements integral over
\(\mathcal O_Z\).  Because \(\mathcal O_Z\) is integrally closed in
its function field, those pullbacks already lie in \(\mathcal O_Z\).
This provides the local ring homomorphisms necessary to lift \(f\)
through \(\widetilde X\to X\).  Their uniqueness follows because
the lift agrees on the common function fields.  Describing this
integrality argument explicitly makes the universal property more
than an abstract arrow diagram.
